% generated from Paper 5/anccft5.tex
\chapter{Algorithmic non-commutative class field theory, V: the degeneracy structure of the Iwahori tower and the Euler correction}\label{chap:V}
\chaptermark{V}
\begin{chapterabstract}
The congruence-model (Iwahori path) tower $Y_0\leftarrow Y_1\leftarrow\cdots$ of
[Ch.~\ref{chap:II}, \S5] is non-normal, carries no deck action, and obeys the exact pseudo-Iwasawa
law $\nu_p=\mu p^k-k+\nu_0$ with $\lambda=-1$, impossible over any commutative
$\Zp[[T]]$ base ([Ch.~\ref{chap:II}, Rem.~5.2], [Ch.~\ref{chap:III}, \S2]). This paper equips that tower with its
canonical algebraic structure and identifies, without fabricating a theory that does
not yet exist, exactly what any future non-commutative characteristic formalism must
specialize to.

The two degeneracy maps --- initial and terminal truncation of non-backtracking paths,
the graph avatars of $X_0(p^{k+1})\rightrightarrows X_0(p^k)$ --- are packaged into an
incidence calculus with six exact identities, among them
$\ta\etp=A_k$ (the level-$k$ Hecke operator realized through level $k{+}1$; at the
bottom, $\ta\etp=T_p$ of [Ch.~\ref{chap:I}]) and $\ta\taup=\et\etp=qI$. Theorem~\ref{V:thm:euler}
computes the degeneracy complex: $\ker(\taup-\etp)=\Z\mathbf1$ and
$\coker(\ta-\et)\cong\Z$ torsion-free, so the Euler rank $\chi=1$ at every level ---
the $-\chi$ of [Ch.~\ref{chap:III}, Thm.~2.4] is the $H^0$ of this complex. Theorem~\ref{V:thm:norm}
shows the tail degeneracy alone intertwines the sandpile Laplacians and the Hecke
operators, is surjective on $p$-primary torsion with kernel of the exact order
$q^{\,n_k(q-1)-1}$ forced by [Ch.~\ref{chap:III}, Thm.~2.1], and hence assembles the tower into one
pro-module $X_\infty=\varprojlim(\coker\Delta_k\otimes\Zp,\ta)$ carrying a limit Hecke
operator; the head degeneracy satisfies the mirror identities for the in-Laplacian
tower and the exact asymmetry $\et\Delta_{k+1}=q(\et-\ta)$, so the structure is
genuinely two-sided. Proposition~\ref{V:prop:nonnormal} records a second, independent
rigorous certificate that no commutative base suffices: $A_kA_k^{\,t}\neq
A_k^{\,t}A_k$ at every computed level. We then state, as axioms rather than as a
theorem, the specialization data any characteristic theory over the Iwahori--Hecke
base must reproduce, and we decline --- for the same reason as in [Ch.~\ref{chap:III}] --- to define
an Akashi-type series whose hypotheses our pro-ring does not satisfy.
\end{chapterabstract}

\section*{Status of the results}

Theorems~\ref{V:thm:euler} and \ref{V:thm:norm} and Propositions~\ref{V:prop:calculus},
\ref{V:prop:mirror} are proved for a general Eulerian $q$-regular strongly connected
level $\ge1$ (the base level $k=0$ is the boundary case, treated in
Remark~\ref{V:rem:level0}); every identity was additionally machine-verified exactly on
the $K_4$ tower at all levels $k\le3$ (Table~\ref{V:tab:battery}).
Proposition~\ref{V:prop:nonnormal} is a finite exact computation on the fixed tower,
hence a proof for it; no general-graph claim is made. Section~\ref{V:sec:axioms} proves
nothing and says so: it is a specification, and Problem~\ref{V:prob:char} is open.

\section{The degeneracy pair and its incidence calculus}\label{V:sec:calculus}

Levels: $Y_0$ a finite connected $(q+1)$-regular graph; for $k\ge1$, $V_{Y_k}$ = the
non-backtracking $k$-paths of $Y_0$, with an arc $p\to p'$ when $p'$ is the shift of a
common $(k{+}1)$-path; $Y_{k+1}=L(Y_k)$ for $k\ge1$, each $Y_k$ Eulerian $q$-regular,
$n_k=n(q+1)q^{k-1}$ ([Ch.~\ref{chap:II}, \S5], [Ch.~\ref{chap:III}, \S2]).

\begin{definition}[Degeneracy maps]\label{V:def:degen}
For $k\ge1$ let $\tau,\eta\colon Y_{k+1}\to Y_k$ be terminal and initial truncation of
paths, $\tau(v_0\cdots v_{k+1})=(v_0\cdots v_k)$, $\eta(v_0\cdots v_{k+1})
=(v_1\cdots v_{k+1})$; under $V_{Y_{k+1}}=\mathrm{Arcs}(Y_k)$ these are the tail and
head maps. Write $\ta,\et\in M_{n_k\times n_{k+1}}(\Z)$ for the pushforwards (sum over
fibres) and $\taup=\ta^{\,t}$, $\etp=\et^{\,t}$ for the pullbacks; thus
$\ta=D_{\mathrm{out}}$, $\et=D_{\mathrm{in}}$, the incidence matrices of
[Ch.~\ref{chap:III}, Thm.~2.1].
\end{definition}

\begin{proposition}[Incidence calculus]\label{V:prop:calculus}
For every $k\ge1$:
\[
\ta\etp=A_k,\quad \et\taup=A_k^{\,t},\quad \etp\ta=A_{k+1},\quad \taup\et=A_{k+1}^{\,t},
\quad \ta\taup=\et\etp=q\,I_{n_k}.
\]
At the boundary level, with $\tau,\eta\colon Y_1\to Y_0$ the endpoint maps on oriented
edges, $\ta\etp=T_p$, the Hecke operator of {\rm[I]}.
\end{proposition}

\begin{proof}
Entrywise bookkeeping of tails and heads; the middle two are [Ch.~\ref{chap:III}, Thm.~2.1, Step~1],
the outer two their transposes, the last two the regularity of degrees. The boundary
identity is $S\,T^{\,t}=T_p$ in the notation of [Ch.~\ref{chap:I}, Thm.~3.7]. All six verified as
exact matrix identities at $k=1,2$ (Table~\ref{V:tab:battery}).
\end{proof}

The pair $(\ta,\et)$ is the graph avatar of the modular degeneracy pair
$X_0(p^{k+1})\rightrightarrows X_0(p^k)$, and Proposition~\ref{V:prop:calculus} is the
statement that the level-$k$ Hecke operator is the correspondence
$\ta\circ\etp$ computed through level $k{+}1$ --- the tower knows its own Hecke
theory.

\section{The degeneracy complex and the Euler rank}\label{V:sec:euler}

\begin{theorem}[Euler rank of the degeneracy complex]\label{V:thm:euler}
For every $k\ge1$, with $\partial_k:=\ta-\et\colon\Z^{V_{Y_{k+1}}}\to\Z^{V_{Y_k}}$ and
$\partial_k^{\,*}:=\taup-\etp$:
\begin{enumerate}
\item[(i)] $\ker\partial_k^{\,*}=\Z\cdot\mathbf1$: the only pulled-back-equal
functions are the constants --- the Perron--Eisenstein line, of rank $\chi_k=1$;
\item[(ii)] $\coker\,\partial_k\cong\Z$ via the degree functional, and it is
torsion-free: the image of $\partial_k$ is exactly the augmentation-zero sublattice.
\end{enumerate}
Hence the Euler rank of the two-term degeneracy complex is $\chi=1$ at every level,
and the correction $-\chi\cdot k=-k$ of {\rm[Ch.~\ref{chap:III}, Thm.~2.4]} is the accumulated
$H^0$ of this complex, once per level.
\end{theorem}

\begin{proof}
(i) $(\partial_k^{\,*}x)_{(a\to b)}=x_a-x_b$; vanishing forces $x$ constant along
arcs, hence constant by strong connectivity. (ii) $\partial_k(e_{(a\to b)})
=e_a-e_b$; over a connected digraph the differences $e_a-e_b$ along arcs generate,
integrally, the full kernel of the degree map (telescope along paths; any
$z$ with $\sum z_i=0$ is an integer combination of consecutive differences), so
$\operatorname{im}\partial_k=\ker(\deg)$ and $\coker\cong\Z$ with no torsion.
Machine confirmation: the Smith form of $\partial_k$ has exactly $n_k-1$ unit
invariants at $k=1,2$ (Table~\ref{V:tab:battery}).
\end{proof}

\section{The tail norm and the pro-module}\label{V:sec:norm}

\begin{theorem}[Tail norm]\label{V:thm:norm}
For every $k\ge1$, with $\Delta_k=qI-A_k$ the out-Laplacian:
\begin{enumerate}
\item[(i)] $\ta\,\Delta_{k+1}=\Delta_k\,\ta$ and $\ta\,A_{k+1}=A_k\,\ta$;
\item[(ii)] $\ta$ is surjective and preserves the degree functional; hence the
induced map on sandpile cokernels is surjective, restricts to a surjection of
$p$-primary torsion, and its kernel there has order exactly
$q^{\,n_k(q-1)-1}$, the increment proved in {\rm[Ch.~\ref{chap:III}, Thm.~2.1]};
\item[(iii)] consequently
$X_\infty:=\varprojlim_k\bigl(\coker_\Z(\Delta_k)\otimes\Zp,\ \ta\bigr)$
is a well-defined pro-module carrying the limit Hecke operator
$A_\infty:=\varprojlim A_k$, and its torsion-size sequence is the pseudo-Iwasawa law
$\ord_p=\mu p^k-k+\nu_0$, $\mu=n(q+1)/q$, of {\rm[Ch.~\ref{chap:III}, Thm.~2.2]} --- now the size
function of a single pro-object rather than a sequence of unrelated groups.
\end{enumerate}
\end{theorem}

\begin{proof}
(i) $\ta A_{k+1}=D_{\mathrm{out}}D_{\mathrm{in}}^{\,t}D_{\mathrm{out}}
=A_kD_{\mathrm{out}}=A_k\ta$ by associativity and
Proposition~\ref{V:prop:calculus}; subtract from $q\ta$. (ii) Every level-$k$ vertex
extends, so $\ta$ hits every basis vector; $\deg\circ\ta=\deg$ since fibres are
summed; torsion-surjectivity then follows exactly as in [Ch.~\ref{chap:IV}, Lem.~1.3] (the kernel of
$\deg$ is the torsion), and the kernel order is forced by counting against
[Ch.~\ref{chap:III}, Thm.~2.1]. (iii) Inverse limits over surjections of compact modules; the Hecke
compatibility is (i). Explicit machine confirmation of (ii) at the transition
$2\to1$: surjective on the $2$-part with kernel order exactly $2^{11}=2^{\,n_1-1}$,
the orders $7,18$ matching [Ch.~\ref{chap:II}, Rem.~5.2] (Table~\ref{V:tab:battery}).
\end{proof}

\begin{proposition}[Head mirror, and the exact asymmetry]\label{V:prop:mirror}
For every $k\ge1$: $\et$ intertwines the in-Laplacians,
$\et\,(qI-A_{k+1}^{\,t})=(qI-A_k^{\,t})\,\et$, so the mirrored construction yields the
in-sandpile pro-module $X_\infty^{\mathrm{in}}$; on the out-tower, however,
\[
\et\,\Delta_{k+1}\;=\;q\,(\et-\ta)\;=\;-q\,\partial_k ,
\]
so $\et$ does not descend to the out-cokernels: the failure is exactly $q$ times the
degeneracy boundary of Theorem~\ref{V:thm:euler}. The tower is thus canonically
two-sided --- a $(\tau,\eta)$-bimodule system, not a single inverse system --- and
the obstruction to collapsing the two sides is again the Euler complex.
\end{proposition}

\begin{proof}
$\et A_{k+1}^{\,t}=D_{\mathrm{in}}D_{\mathrm{out}}^{\,t}D_{\mathrm{in}}
=A_k^{\,t}D_{\mathrm{in}}$; and
$\et\Delta_{k+1}=qD_{\mathrm{in}}-D_{\mathrm{in}}D_{\mathrm{in}}^{\,t}D_{\mathrm{out}}
=q(D_{\mathrm{in}}-D_{\mathrm{out}})$ by
$D_{\mathrm{in}}D_{\mathrm{in}}^{\,t}=qI$. Verified at $k=1,2$.
\end{proof}

\begin{remark}[The boundary level]\label{V:rem:level0}
At $k=0$ the base is undirected and $(q+1)$-regular, and the clean intertwining fails
by exactly the degeneracy boundary:
$\ta\Delta_1-\Delta_0\ta=-(\ta-\et)=-\partial_0$, verified as an exact identity. The
Eisenstein special layer of the whole series --- the $(q+1)$-versus-$q$ mismatch of
[Ch.~\ref{chap:I}, Rem.~1.8], the shared Perron root of [Ch.~\ref{chap:II}], the excluded trivial eigenvalue of
[Ch.~\ref{chap:III}] --- reappears here as the unique level where the norm formalism acquires a
correction term, and the correction is once more $\partial$.
\end{remark}

\begin{proposition}[Second non-commutativity certificate]\label{V:prop:nonnormal}
On the $K_4$ tower, $A_kA_k^{\,t}\neq A_k^{\,t}A_k$ for $k=1,2,3$ (exact
computation, hence a proof for this tower). Together with the growth-theoretic
certificate $\lambda=-1<0$ of {\rm[Ch.~\ref{chap:III}, Prop.~2.3]}, this gives two independent
rigorous witnesses that the level operator algebras
$\Z\langle A_k,A_k^{\,t}\rangle$, and a fortiori any base ring over which the tower's
characteristic theory could live, are non-commutative: the first says the
\emph{growth} cannot come from $\Zp[[T]]$, the second that the \emph{operators}
themselves do not commute with their adjoints.
\end{proposition}

\section{What a characteristic theory must specialize to}\label{V:sec:axioms}

We now state --- as a specification, not a theorem --- the data pinned down by
Papers~II--V that any characteristic formalism $\mathrm{Char}$ for
$(\tau,\eta)$-towers over the (pro-)Iwahori--Hecke base must reproduce:
\begin{enumerate}
\item[(A1)] a multiplicative rank-density term recovering $\mu=n(q+1)/q$, the arc
proliferation of {\rm[Ch.~\ref{chap:III}, Thm.~2.2]};
\item[(A2)] an additive Euler term $-\chi\cdot k$ with
$\chi=\rk\ker(\taup-\etp)$, the $H^0$ of the degeneracy complex
(Theorem~\ref{V:thm:euler}) --- intrinsic to the base, external to any commutative
characteristic ideal;
\item[(A3)] on sub-towers carrying an abelian deck action, reduction to the
commutative characteristic ideal $\bigl(\det D(1+T)\bigr)$ of {\rm[Ch.~\ref{chap:IV}, Thm.~2.1]},
with the Newton-polygon evaluation of {\rm[Ch.~\ref{chap:II}, Rem.~5.4]}.
\end{enumerate}

\begin{problem}\label{V:prob:char}
Construct $\mathrm{Char}$ satisfying (A1)--(A3) on the pro-module
$(X_\infty,A_\infty,\tau,\eta)$ of Theorem~\ref{V:thm:norm}.
\end{problem}

We explicitly decline to define an Akashi-type series here. The non-commutative
Iwasawa formalism in which Akashi series live is built over Iwasawa algebras of
compact $p$-adic Lie groups \cite{CFKSV}; our pro-ring is the inverse limit of the
finite-dimensional, non-normal level algebras of
Proposition~\ref{V:prop:nonnormal}, generated by degeneracy correspondences rather
than by a group, and none of the hypotheses of that formalism are available. Writing
down an ``Akashi series'' whose only property is to output the already-proved growth
law would be circular dressing, and this series does not do that
([Ch.~\ref{chap:III}, Rem.~2.6]). What Papers~II--V have done instead is to make
Problem~\ref{V:prob:char} well-posed: the category, the operators, the two invariants
any answer must produce, and the commutative shadow it must reduce to, are now all
theorems.

\section{The verification battery}\label{V:sec:battery}

\begin{table}[ht]
\centering\small
\renewcommand{\arraystretch}{1.2}
\begin{tabular}{lcl}
\toprule
check & levels & result\\
\midrule
(F) six incidence identities, incl.\ $\ta\etp=A_k$, $\ta\taup=\et\etp=qI$ & $k=1,2$ & exact\\
(T) $\ta\Delta_{k+1}=\Delta_k\ta$,\ $\ta A_{k+1}=A_k\ta$ & $k=1,2$ & exact\\
(H) in-mirror for $\et$; asymmetry $\et\Delta_{k+1}=q(\et-\ta)$ & $k=1,2$ & exact\\
(E) $\ker(\taup-\etp)=\Z\mathbf1$; $\mathrm{SNF}(\ta-\et)$: $n_k{-}1$ units & $k=1,2$ & exact\\
(N) tail norm on $2$-part: surjective, $|\ker|=2^{11}=2^{\,n_1-1}$ & $2\to1$ & exact\\
(C) $A_kA_k^{\,t}\neq A_k^{\,t}A_k$ & $k=1,2,3$ & exact\\
(Z) $\ta\Delta_1-\Delta_0\ta=-\partial_0$;\ \ $\ta\etp=T_p$ at the bottom & $0\to1$ & exact\\
\bottomrule
\end{tabular}
\medskip
\caption{Exact integer verification (scripts \texttt{iw\_stage1--3.py}) on the $K_4$
tower, $q=2$, levels $\le3$ (up to $48$ vertices). Each general lemma is proved in
the text; for the fixed tower each line is itself a proof.}
\label{V:tab:battery}
\end{table}

\section*{Acknowledgement of provenance and caveats}

Unrefereed preprint. All numbered statements are proved, with
Proposition~\ref{V:prop:nonnormal} proved for the fixed $K_4$ tower by exact
computation and asserted for no other graph. Section~\ref{V:sec:axioms} is a
specification; Problem~\ref{V:prob:char} is open, and the paper's deliberate refusal to
define an Akashi-type invariant is explained there. The single external citation
\cite{CFKSV} is contextual and load-bearing for nothing; its bibliographic data are
now primary-source verified.
